\documentclass{article}
\usepackage{amssymb}
\usepackage{amsmath}
\usepackage{xcolor}

\title{Homework 7}
\author{Brandon Reich}
\date{6 March 2026}

\begin{document}
\maketitle

\subsection*{5.1.1}
\begin{itemize}
    \item a) $\{(1,2),\ (2,1),\ (3,3)\}$
\end{itemize}

\subsection*{5.1.4}
\begin{itemize}
    \item a) Domain of $C$ is $\{0,1\}^3$. $xCy$ if $y$ is obtained from $x$ by changing one 0 to a 1.\\
    Pairs: $(000,001),\ (000,010),\ (000,100),\ (001,011),\ (001,101),\ (010,011),\ (010,110),\ (011,111),\ (100,101),\ (100,110),\ (101,111),\ (110,111)$
\end{itemize}

\subsection*{5.2.1}
\begin{itemize}
    \item b) $E$ on $\mathbb{R}$, $xEy$ if $x \le y$.\\
    \textbf{Reflexive}: Yes --- $x \le x$ for all $x$.\\
    \textbf{Anti-symmetric}: Yes --- if $x \le y$ and $y \le x$, then $x = y$.\\
    \textbf{Not symmetric}: $1 \le 2$ but $2 \not\le 1$.\\
    \textbf{Transitive}: Yes --- if $x \le y$ and $y \le z$, then $x \le z$.
\end{itemize}

\subsection*{5.2.4}
\begin{itemize}
    \item b) $R = \{(a,b),\ (b,c),\ (a,c)\}$ on $\{a,b,c\}$.\\
    \textbf{Anti-reflexive}: Yes --- no pair $(x,x)$ is in $R$.\\
    \textbf{Anti-symmetric}: Yes --- no pair $(x,y)$ and $(y,x)$ both in $R$ with $x \ne y$.\\
    \textbf{Transitive}: Yes --- $(a,b)$ and $(b,c) \Rightarrow (a,c) \in R$. All other required pairs check out.
\end{itemize}

\subsection*{5.3.1}
\begin{itemize}
    \item a) In-degree of $d$ = 2 (edges $b \to d$ and $c \to d$).
    \item b) Out-degree of $c$ = 1 (edge $c \to d$).
    \item c) The head of edge $(b, c)$ is $c$.
    \item f) $\langle a, g, f, c, d \rangle$: Yes, it is a walk ($a \to g,\ g \to f,\ f \to c,\ c \to d$ all exist). It is an open walk since $a \ne d$. No edge is repeated, so it is a trail. No vertex is repeated, so it is a path.
    \item h) $\langle c, g, f, e \rangle$: It is not a circuit because $c \ne e$ (the walk is not closed). Therefore it is not a cycle either.
\end{itemize}

\subsection*{5.3.4}
\begin{itemize}
    \item b) Paths exist of length $x = 0, 1, 2$.\\
    Length 0: trivial path at any single vertex.\\
    Length 1: e.g., $1 \to 2$.\\
    Length 2: e.g., $1 \to 2 \to 3$.\\
    Length 3 is impossible --- would require 4 distinct vertices, but only 3 exist.
    \item d) Cycles exist of length $x = 1, 2$.\\
    Length 1: self-loop $2 \to 2$.\\
    Length 2: e.g., $1 \to 2 \to 1$ or $2 \to 3 \to 2$.\\
    Length 3 is impossible --- would need a cycle through all 3 vertices, but there is no edge $3 \to 1$ or $1 \to 3$.
\end{itemize}

\subsection*{5.4.1}
$S = \{(a,b),\ (a,c),\ (c,d),\ (c,a)\}$, \quad $R = \{(b,c),\ (c,b),\ (a,d),\ (d,b)\}$
\begin{itemize}
    \item a) $S \circ R$: For each $(x,y) \in R$, find $(y,z) \in S$:\\
    $(b,c) \in R$: $(c,d),\ (c,a) \in S \Rightarrow (b,d),\ (b,a)$\\
    $(c,b) \in R$: $(b,?) \in S$: none\\
    $(a,d) \in R$: $(d,?) \in S$: none\\
    $(d,b) \in R$: $(b,?) \in S$: none\\
    $S \circ R = \{(b,a),\ (b,d)\}$

    \item c) $S \circ S$: For each $(x,y) \in S$, find $(y,z) \in S$:\\
    $(a,b)$: $(b,?) \in S$: none\\
    $(a,c)$: $(c,d),\ (c,a) \in S \Rightarrow (a,d),\ (a,a)$\\
    $(c,d)$: $(d,?) \in S$: none\\
    $(c,a)$: $(a,b),\ (a,c) \in S \Rightarrow (c,b),\ (c,c)$\\
    $S \circ S = \{(a,a),\ (a,d),\ (c,b),\ (c,c)\}$
\end{itemize}

\subsection*{5.4.3}
$A = \{(a,a),\ (a,c)\}$, \quad $B = \{(a,a),\ (c,d)\}$
\begin{itemize}
    \item a) $B \circ A$: For each $(x,y) \in A$, find $(y,z) \in B$:\\
    $(a,a) \in A$: $(a,a) \in B \Rightarrow (a,a)$\\
    $(a,c) \in A$: $(c,d) \in B \Rightarrow (a,d)$\\
    $B \circ A = \{(a,a),\ (a,d)\}$
\end{itemize}

\subsection*{5.5.1}
Graph $G$ edges: $a \to b,\ a \to g,\ b \to b,\ b \to c,\ b \to d,\ c \to d,\ g \to c,\ g \to f,\ f \to c,\ f \to e$
\begin{itemize}
    \item a) Is $(a,b)$ in $G^2$? Yes --- $a \to b \to b$ (using the self-loop on $b$).
    \item c) Is $(g,g)$ in $G^3$? No --- the only edge into $g$ is $a \to g$, and no edge leads to $a$, so there is no walk of any length from $g$ back to $g$.
\end{itemize}

\subsection*{5.5.2}
$G = \{(1,2),\ (2,1),\ (2,3),\ (4,2),\ (4,3)\}$
\begin{itemize}
    \item a) $G^2 = \{(1,1),\ (1,3),\ (2,2),\ (4,1),\ (4,3)\}$\\
    $G^3 = \{(1,2),\ (2,1),\ (2,3),\ (4,2)\}$\\
    $G^4 = \{(1,1),\ (1,3),\ (2,2),\ (4,1),\ (4,3)\}$\\
    $G^+ = G \cup G^2 \cup G^3 \cup G^4 = \{(1,1),\ (1,2),\ (1,3),\ (2,1),\ (2,2),\ (2,3),\ (4,1),\ (4,2),\ (4,3)\}$
\end{itemize}

\subsection*{5.5.3}
\begin{itemize}
    \item a) Same graph as 5.5.2. The transitive closure is:\\
    $G^+ = \{(1,1),\ (1,2),\ (1,3),\ (2,1),\ (2,2),\ (2,3),\ (4,1),\ (4,2),\ (4,3)\}$
\end{itemize}

\subsection*{5.5.4}
Given $G^2$ and $G^3$ for graph $G$ on $\{1,2,3,4,5\}$.
\begin{itemize}
    \item c) Is there any closed walk of length 3 in $G$?\\
    Yes --- a closed walk of length 3 exists if $(x,x) \in G^3$ for some $x$. Vertices 1, 2, 3, and 5 all have self-loops in $G^3$.

    \item e) Is there a walk of length 5 from vertex 2 to vertex 3 in $G$?\\
    Yes --- $G^5 = G^2 \cdot G^3$. From $G^2$: $(2,3) \in G^2$. From $G^3$: $(3,3) \in G^3$. So $(2,3) \in G^5$.
\end{itemize}

\subsection*{5.6.1}
$G = \{(1,2),\ (2,1),\ (2,3),\ (4,2),\ (4,3)\}$
\begin{itemize}
    \item a) Adjacency matrix $A$ and powers:\\
    $A = \begin{bmatrix} 0 & 1 & 0 & 0 \\ 1 & 0 & 1 & 0 \\ 0 & 0 & 0 & 0 \\ 0 & 1 & 1 & 0 \end{bmatrix}$\quad
    $A^2 = \begin{bmatrix} 1 & 0 & 1 & 0 \\ 0 & 1 & 0 & 0 \\ 0 & 0 & 0 & 0 \\ 1 & 0 & 1 & 0 \end{bmatrix}$\\[6pt]
    $A^3 = \begin{bmatrix} 0 & 1 & 0 & 0 \\ 1 & 0 & 1 & 0 \\ 0 & 0 & 0 & 0 \\ 0 & 1 & 0 & 0 \end{bmatrix}$\quad
    $A^4 = \begin{bmatrix} 1 & 0 & 1 & 0 \\ 0 & 1 & 0 & 0 \\ 0 & 0 & 0 & 0 \\ 1 & 0 & 1 & 0 \end{bmatrix}$\\[6pt]
    $G^+ = A \lor A^2 \lor A^3 \lor A^4 = \begin{bmatrix} 1 & 1 & 1 & 0 \\ 1 & 1 & 1 & 0 \\ 0 & 0 & 0 & 0 \\ 1 & 1 & 1 & 0 \end{bmatrix}$
\end{itemize}

\subsection*{5.6.3}
$A^2 = \begin{bmatrix} 0&1&0&0&0 \\ 0&0&1&0&0 \\ 1&0&0&0&0 \\ 1&0&0&1&0 \\ 0&1&1&0&1 \end{bmatrix}$\quad
$A^3 = \begin{bmatrix} 1&0&0&0&0 \\ 0&1&0&0&0 \\ 0&0&1&0&0 \\ 0&1&1&0&1 \\ 1&1&0&1&0 \end{bmatrix}$
\begin{itemize}
    \item a) Which vertices can reach vertex 2 by a walk of length 3?\\
    Look at column 2 of $A^3$: entries are 0, 1, 0, 1, 1.\\
    Vertices $\{2, 4, 5\}$ can reach vertex 2 by a walk of length 3.

    \item d) Is $(2,2)$ in the transitive closure of $G$?\\
    Yes --- $A^3[2,2] = 1$, so there is a walk of length 3 from 2 to 2. Since $(2,2) \in G^3$, it is in $G^+$.
\end{itemize}

\subsection*{5.7.1}
Hasse diagram on $\{A,B,C,D,E,F,G,H,I,J\}$:\\
$A < B < C < D$, $B < E < D$, $B < E < G$, $I < H$, $J < H$, $F$ is isolated.
\begin{itemize}
    \item a) Minimal elements: $A,\ F,\ I,\ J$ (no element below them).
    \item b) Maximal elements: $D,\ G,\ H,\ F$ (no element above them).
    \item c) Comparable pairs from the list:\\
    $(A, D)$: Yes --- $A < B < C < D$.\\
    $(J, F)$: No --- incomparable (different components).\\
    $(B, E)$: Yes --- $B < E$.\\
    $(G, F)$: No --- incomparable.\\
    $(D, B)$: Yes --- $B < C < D$.\\
    $(C, F)$: No --- incomparable.\\
    $(H, I)$: Yes --- $I < H$.\\
    $(C, E)$: No --- incomparable (both above $B$, both below $D$, but no path between them).
\end{itemize}

\subsection*{5.7.2}
\begin{itemize}
    \item a) $P(A)$ where $A = \{x, y\}$. $X \preceq Y$ if $X \subseteq Y$.\\
    $P(A) = \{\emptyset,\ \{x\},\ \{y\},\ \{x,y\}\}$.\\
    Hasse diagram (bottom to top):\\
    $\emptyset$ at bottom, $\{x\}$ and $\{y\}$ in the middle (incomparable), $\{x,y\}$ at top.\\
    Edges: $\emptyset - \{x\}$, $\emptyset - \{y\}$, $\{x\} - \{x,y\}$, $\{y\} - \{x,y\}$.
\end{itemize}

\subsection*{5.8.1}
\begin{itemize}
    \item a) Domain: all English words. $x$ related to $y$ if $x$ appears before $y$ alphabetically.\\
    \textbf{Anti-reflexive}: Yes --- a word does not appear before itself.\\
    \textbf{Anti-symmetric}: Yes --- if $x$ before $y$, then $y$ is not before $x$.\\
    \textbf{Transitive}: Yes --- if $x$ before $y$ and $y$ before $z$, then $x$ before $z$.\\
    This is a \textbf{strict order}. It is also a \textbf{strict total order} because for any two distinct words, one appears before the other.
\end{itemize}

\subsection*{5.8.2}
DAG edges: $b \to c,\ b \to d,\ c \to a,\ f \to e,\ e \to a,\ a \to g$.
\begin{itemize}
    \item a) Two topological sorts:\\
    1. $b,\ f,\ d,\ c,\ e,\ a,\ g$\\
    2. $f,\ b,\ c,\ d,\ e,\ a,\ g$
\end{itemize}

\subsection*{5.8.4}
\begin{itemize}
    \item a) Edges: $d \to a,\ d \to b,\ d \to e,\ a \to b,\ b \to e,\ c \to a,\ c \to b$.\\
    \textbf{Acyclic?} Yes --- no vertex can reach itself; all paths lead to $e$ (a sink).\\[4pt]
    \textbf{Transitive closure:}\\
    From $d$: $d \to a,\ d \to b,\ d \to e$ (and $a \to b \to e$, so already covered).\\
    From $a$: $a \to b,\ a \to e$ (via $b \to e$).\\
    From $b$: $b \to e$.\\
    From $c$: $c \to a,\ c \to b,\ c \to e$ (via $a \to b \to e$).\\
    From $e$: nothing.\\[4pt]
    $G^+ = \{(d,a),\ (d,b),\ (d,e),\ (a,b),\ (a,e),\ (b,e),\ (c,a),\ (c,b),\ (c,e)\}$
\end{itemize}

\subsection*{5.9.1}
\begin{itemize}
    \item b) Domain is a group of people. $xMy$ if $x$ and $y$ have the same biological mother.\\
    \textbf{Reflexive}: Yes --- every person has the same biological mother as themselves.\\
    \textbf{Symmetric}: Yes --- if $x$ and $y$ share the same mother, then $y$ and $x$ do too.\\
    \textbf{Transitive}: Yes --- if $x$ and $y$ share a mother, and $y$ and $z$ share a mother, then $x$ and $z$ share that same mother (each person has exactly one biological mother).\\[4pt]
    $M$ \textbf{is an equivalence relation}. The equivalence classes are sets of people who have the same biological mother (i.e., maternal siblings).
\end{itemize}

\subsection*{5.9.2}
\begin{itemize}
    \item a) $S = \{7,2,13,44,56,34,99,31,4,17\}$. $xDy$ if $x \bmod 4 = y \bmod 4$.\\[4pt]
    Remainders:\\
    $7 \bmod 4 = 3$, \quad $2 \bmod 4 = 2$, \quad $13 \bmod 4 = 1$, \quad $44 \bmod 4 = 0$\\
    $56 \bmod 4 = 0$, \quad $34 \bmod 4 = 2$, \quad $99 \bmod 4 = 3$, \quad $31 \bmod 4 = 3$\\
    $4 \bmod 4 = 0$, \quad $17 \bmod 4 = 1$\\[4pt]
    Partition by equivalence classes:\\
    $[44] = \{44,\ 56,\ 4\}$ \quad (remainder 0)\\
    $[13] = \{13,\ 17\}$ \quad (remainder 1)\\
    $[2] = \{2,\ 34\}$ \quad (remainder 2)\\
    $[7] = \{7,\ 99,\ 31\}$ \quad (remainder 3)
\end{itemize}

\end{document}
